% §3 Related Work
% Status: Written 2026-03-22 — based on pre-experiment draft, updated with E1–E7 findings.

\section{Related Work}
\label{sec:related}

\subsection{Exascale Programming in Practice}

Evans et al.~\cite{Evans2022} present an empirical survey of 62 application codes developed
under the Exascale Computing Project, examining the programming languages and
programming models adopted in practice for emerging exascale systems. The study
shows that exascale applications depend heavily on established programming
languages, with C++ becoming dominant (used in $\sim$85\% of codes) and Fortran
in rapid decline, and employ hybrid programming models combining MPI for
internode parallelism with GPU-centric approaches such as CUDA, HIP, OpenMP,
and portability frameworks like Kokkos and RAJA. Rather than adopting
fundamentally new programming paradigms, application teams achieve productivity
through layered abstractions, co-design middleware (e.g., AMReX, CEED, CoPA),
and selective use of native implementations when portability layers prove
insufficient. Notably, $\sim$55\% of codes still use native or custom
implementations. The survey reveals that productivity at exascale is achieved
not by fully hiding parallelism or hardware details, but by enabling developers
to choose where abstraction applies and where explicit control must be
retained---explicit parallelism, platform-specific tuning, and model selection
remain unavoidable in practice.

Heroux et al.~\cite{heroux2020ecp} present an ecosystem-level overview of the
software libraries and tools developed within the Exascale Computing Project
(ECP), documenting how programming models, runtimes, and supporting
infrastructure were designed and deployed to enable productivity on emerging
exascale systems. Rather than proposing new programming languages or paradigms,
the paper shows that exascale productivity is pursued through a layered software
stack in which portability frameworks (e.g., Kokkos, RAJA), runtimes, and
co-designed libraries collectively absorb hardware heterogeneity across AMD,
Intel, and NVIDIA platforms. Programming models are positioned as infrastructure
components that support libraries and applications, with no convergence toward a
single abstraction; instead, multiple approaches coexist to accommodate diverse
application domains, legacy codebases, and vendor-specific constraints. The
paper emphasizes that effective productivity at exascale depends on ecosystem
interoperability, vendor co-design, and standardized software delivery
mechanisms (e.g., Spack, E4S), while acknowledging that expert involvement,
continuous tuning, and architecture-specific optimizations remain unavoidable.

\textit{Together, these studies establish that exascale programming model design
is shaped less by the search for a universal abstraction and more by the
deliberate construction of interoperable software ecosystems that allow selective
abstraction while preserving expert control---a finding that directly motivates
the need for a principled, measurement-driven framework for abstraction
selection.}

% ------------------------------------------------------------------
\subsection{Task-Based and Dataflow Abstractions}
% ------------------------------------------------------------------

Bon\'e et al.~\cite{bone2024taskbased} propose a task-based data-flow
methodology that enables the coordinated use of multiple accelerator programming
models---CUDA, SYCL, Triton, and OpenMP offload---within a single application.
The work addresses the growing challenge of API multiplicity and runtime
interference in heterogeneous systems, where independent runtimes contend for
CPU resources and introduce synchronization complexity. By expressing
applications as task DAGs and using task-aware libraries to defer dependency
release until device operations complete, the approach shifts synchronization
and thread management into the runtime, mitigating oversubscription and reducing
manual coordination. However, developers still write native kernels, annotate
dependencies, and tune task granularity, positioning task-based dataflow as a
coordination abstraction that reorganizes complexity rather than eliminating it.

Bauer et al.~\cite{bauer2012legion} investigate the limits of high-productivity
programming systems based on fully implicit parallelism by addressing dynamic
dependence analysis and distributed coherence in distributed-memory
environments. The paper demonstrates that apparently sequential semantics can be
preserved even in the presence of data-dependent control flow and overlapping,
aliased data partitions by shifting all responsibility for dependence discovery,
synchronization, and data movement into the runtime system. By reducing
coherence and dependence analysis to the visibility problem from computer
graphics and implementing scalable ray-casting algorithms within the Legion
runtime, the authors show that content-based coherence is technically feasible
without restricting program expressivity. However, achieving this level of
abstraction requires sophisticated distributed metadata structures, non-trivial
runtime overhead, coarse-grained tasks, and substantial engineering effort,
establishing fully implicit parallelism as an upper bound on abstraction rather
than a practical strategy for contemporary exascale programming systems.

John et al.~\cite{john2023multgpu} enhance the PaRSEC task-based dataflow
runtime with dynamic inter-GPU work sharing to address post-mapping load
imbalance in multi-GPU nodes, a problem acute for irregular applications where
static heuristics fail. By enabling busy GPUs to push tasks to starving GPUs
via a work-sharing protocol, the runtime absorbs GPU selection, data movement,
and synchronization without programmer intervention, demonstrating 3--12\%
speedup on real sparse matrix workloads and stronger gains (up to $2\times$ at
8 GPUs) on synthetic benchmarks with high imbalance. However, the approach
fails to resolve fundamental scalability limits: migration increases memory
thrashing, introduces lock contention during task selection, and relies on
task-count-based starvation detection that ignores execution time variability
and critical path structure. Ultimately, John et al.\ show that dynamic work
sharing is a necessary but insufficient mechanism for exascale multi-GPU
programming: it mitigates imbalance but does not eliminate the need for careful
task decomposition, data partitioning, and runtime parameter tuning.

\textit{Collectively, these works demonstrate that task-based and dataflow
abstractions can hide coordination complexity and improve productivity, but
consistently displace rather than eliminate the need for architecture-aware
tuning---the granularity, decomposition, and scheduling decisions that determine
whether abstraction delivers its promised performance gains remain the
programmer's responsibility.}

% ------------------------------------------------------------------
\subsection{Adaptive and Self-Aware Runtimes}
% ------------------------------------------------------------------

Chandrasekar et al.~\cite{chandrasekar2023adaptive} investigate adaptive runtime
techniques for node-aware resource optimization in the Charm++ runtime, focusing
on dynamically adjusting core utilization, communication parallelism, message
aggregation strategies, and load balancing based on observed runtime behavior.
The work demonstrates that reactive, runtime-driven adaptation can yield
significant performance and energy benefits for certain workloads, including up
to 40\% speedup for memory-bound stencil codes through core throttling and
multi-fold improvements for communication-heavy applications via adaptive
aggregation and diffusion-based load balancing. However, the study also shows
that adaptivity introduces substantial new complexity in the form of profiling
overheads, runtime exploration phases, numerous exposed tuning parameters, and
mandatory programmer instrumentation. Overall, the paper positions runtime
adaptivity as a mechanism that reorganizes tuning effort rather than eliminating
it, providing measurable gains for well-understood workloads while remaining
sensitive to application characteristics and offering limited relief from
expert-driven performance engineering.

Thomadakis and Chrisochoides~\cite{thomadakis2023prema} present PREMA, a runtime
system that targets performance portability on heterogeneous distributed
platforms by shifting scheduling, data movement, and coherence management from
the programmer to the runtime. The approach introduces task- and data-centric
abstractions (\texttt{hetero\_task}, \texttt{hetero\_object}) that enable
applications to execute across CPUs and GPUs without explicit memory transfers
or synchronization, while dynamically selecting devices and overlapping
communication with computation. Experimental results demonstrate that PREMA can
achieve performance comparable to, and in some cases exceeding, hand-tuned
MPI+CUDA implementations for selected workloads, particularly when aggressive
over-decomposition and runtime-level optimizations are employed. However, truly
portable performance is not fully automatic: programmers must still write
device-specific kernels in a non-standard dialect, manually configure thread
dimensions, and tune task decomposition strategies. As a result, PREMA
effectively automates orchestration and coordination concerns, but it does not
eliminate the need for architecture-aware optimization effort at the application
level.

\textit{These systems illustrate the ceiling of runtime-driven abstraction: even
the most sophisticated adaptive and self-aware runtimes achieve portable
performance only for workloads whose characteristics align with the runtime's
adaptation model, and require significant programmer effort to configure,
instrument, and tune when that alignment breaks down.}

% ------------------------------------------------------------------
\subsection{Performance Portability Layers}
% ------------------------------------------------------------------

Pennycook et al.~\cite{Pennycook2019} introduce the Performance Portability
Coefficient (PPC), a formal metric for quantifying how consistently an
application delivers efficiency across a set of hardware platforms. By
formulating portability as the harmonic mean of per-platform efficiencies
relative to best-known implementations, the metric captures the intuition that
a portability claim is only as strong as its weakest platform. Our study
operationalizes this metric as the primary evaluation instrument across all
190+ measured configurations; the three-tier threshold system we adopt
(PPC $\geq 0.80$ excellent; $0.60$--$0.80$ acceptable; $< 0.60$ poor) is
consistent with Pennycook et al.'s recommended interpretation guidelines.

Zhao et al.~\cite{zhao2023hiplz} introduce HIPLZ, a compiler and runtime system
that enables performance portability of HIP applications across GPU vendors by
implementing the HIP runtime API on top of Intel's Level Zero driver. Rather
than proposing a new programming model, HIPLZ preserves existing HIP semantics
and toolchains, translating HIP kernels to SPIR-V and executing them on Intel
GPUs without source-level changes. Evaluation across 50 HIP applications shows
high functional portability and performance parity with OpenCL on Intel hardware,
demonstrating that narrow, translation-based approaches can reduce vendor lock-in
with minimal adoption friction. However, incomplete API coverage, correctness
failures, and reliance on driver-specific behaviors reveal that thin translation
layers shift low-level complexity into the runtime rather than eliminating it,
limiting robustness and reinforcing the need for native implementations in
performance-critical exascale codes.

Godoy et al.~\cite{Godoy2023} evaluate the performance portability of Julia,
Python/Numba, and Kokkos against vendor-native C/OpenMP, CUDA, and HIP
implementations on exascale-class AMD EPYC/MI250X (CDNA2) and Arm
Altra/NVIDIA A100 nodes using a dense matrix multiplication kernel. On AMD
hardware, Julia achieves 0.903--0.976 efficiency relative to vendor baselines,
with AMDGPU.jl occasionally outperforming HIP; Kokkos achieves 0.677--1.014
efficiency on CPUs but collapses to 0.208--0.260 on NVIDIA A100 GPUs;
Python/Numba consistently underperforms (0.095--0.713) and lacks AMD GPU
support. The study demonstrates that high-level abstractions do not guarantee
performance portability: Julia succeeds on CDNA2 hardware but incurs a 40\%
penalty on NVIDIA A100 single precision, while Kokkos and Python/Numba fail
catastrophically on NVIDIA GPUs despite using vendor backends. Our experiments
extend this finding to CDNA3 (MI300X): the Julia column-major advantage observed
by Godoy et al.\ on CDNA2 (MI250X) is absent on CDNA3, where Julia/ROCm
delivers only $0.31\times$ native DGEMM throughput---a platform-generation
reversal that underscores the architecture-specificity of abstraction outcomes
and motivates evaluation across successive hardware generations.

Deakin et al.~\cite{Deakin2019} introduce GPU-STREAM (BabelStream), a
multi-framework memory bandwidth benchmark that implements the STREAM Triad
pattern across CUDA, OpenCL, OpenMP, Kokkos, RAJA, SYCL, and other frameworks
on a common hardware target. By isolating memory bandwidth from algorithmic
complexity, BabelStream reveals that even a trivially portable workload
exposes implementation-quality differences across frameworks: framework overhead,
memory allocation strategies, and compiler code generation all affect achieved
bandwidth. We use BabelStream as our E1 baseline (STREAM Triad), establishing
the performance ceiling for memory-bound portability before moving to
computationally complex workloads.

Ziogas et al.~\cite{ziogas2025dace} present Data-Centric Python, a
high-performance subset of Python built on the Stateful Dataflow multiGraph
(SDFG) intermediate representation, which enables automatic optimization and
specialization of annotated Python code across CPUs, GPUs, and FPGAs. By
making data movement explicit through a data-centric IR and applying multi-level
dataflow transformations, the framework achieves 2.47$\times$ and 3.75$\times$
speedups over prior best solutions on CPU and GPU respectively, and scales to
512 nodes with up to 93.16\% efficiency on the Piz Daint supercomputer.
Crucially, the system provides a selective control mechanism: power users can
override automatic decisions and specify explicit parallelism and partitioning
schemes when the automatic heuristics are insufficient, while domain scientists
retain Python's productivity for the common case. However, the framework's
performance depends on the quality of its auto-optimization heuristics, which
can fail silently on irregular control flow, non-standard memory access
patterns, or architectures not covered by the backend specialization
pipeline---precisely the conditions under which the abstraction boundary must
be manually managed.

\textit{Across these portability layers, a consistent pattern emerges:
abstraction delivers portable performance for well-structured, regular workloads
on mature backends, but degrades unpredictably when workload characteristics,
compiler heuristics, or hardware targets fall outside the abstraction's design
envelope---and no existing framework provides a principled method for predicting
or detecting these failure modes in advance.}

% ------------------------------------------------------------------
\subsection*{Gap Statement}
% ------------------------------------------------------------------

Despite extensive work on high-level programming models, task-based runtimes,
adaptive execution systems, and portability layers for exascale computing, there
is no systematic, measurement-driven methodology for deciding where abstraction
should be applied and where explicit control must be retained. Existing studies
either (i) propose new abstractions or runtimes and evaluate them in isolation,
or (ii) measure performance portability within a single abstraction layer,
without analyzing cross-layer failure patterns, their root causes, or their
dependence on problem size and hardware generation. As a result, abstraction
choices remain heuristic, experience-driven, and opaque, forcing developers to
rely on trial-and-error rather than evidence when balancing productivity against
performance. What is missing is an empirical, cross-layer characterization of
abstraction effectiveness that explains when abstraction helps, when it fails,
and why---across workload types, problem sizes, and hardware architectures. The
methodology proposed in Section~\ref{sec:methodology} is designed to fill
precisely this gap, providing a measurement-driven decision framework grounded
in systematic empirical evidence from seven representative kernels across two
GPU architectures.
